\chapter{Knowledge conceptualization}

In this chapter we present an analysis of the process that led to expressing the
domain knowledge in logical form. In particular, we discuss the objects that were
identified, together with the relations and properties of those objects, and we
give a conceptual overview of the two complementary models through which the
system reasons under uncertainty. The detailed calculus that underpins the two
models---the certainty-factor combination, the trapezoidal membership functions,
the centroid defuzzification and the weighted-mean scoring---is deferred to the
Implementation chapter, which owns the equations and their derivations; here we
remain at the conceptual level.

The conceptualization is organised around the following steps, which structure
the reasoning from the elementary dialogue with the user up to the final
recommendation:

\begin{enumerate}[label=(\Roman*)]
    \item \emph{Domain objects} ---the beer styles and the named beers the system
    can recommend, each carrying the symbolic descriptors (colour, alcohol,
    flavour, fermentation, carbonation) over which the qualitative reasoning is
    performed, encoded by the \code{beer} template;
    \item \emph{Data-grounded descriptions} ---the authoritative per-style sensory
    profiles, numeric ranges and food pairings extracted from the CraftBeer~2017
    Beer Styles Study Guide, encoded by the \code{beer-profile},
    \code{beer-range} and \code{pairing} templates;
    \item \emph{Attributes and evidence} ---the graded pieces of evidence about
    the user's ideal profile, accumulated on the colour, alcohol, flavour,
    fermentation and carbonation \emph{axes}, encoded by the \code{attribute}
    template in the canonical \emph{best-*} form;
    \item \emph{The dialogue / evidence model} ---the conversational machinery
    that elicits the user's desires through scenario, preference and meal
    questions, tracks the conversation and overlays localised labels, encoded by
    the \code{UI-state}, \code{state-list}, \code{translation} and
    \code{meal-keyword} templates;
    \item \emph{The two uncertainty models} ---the MYCIN certainty-factor model
    and the data-grounded fuzzy-logic model, the latter introducing linguistic
    variables, desires and direct suggestions, encoded by the \code{ling-var},
    \code{ling-term}, \code{desire}, \code{desire-cat}, \code{suggest-name} and
    \code{suggest-style} templates;
    \item \emph{The output} ---the ranked beer recommendations and their
    data-grounded rationale, encoded by the \code{attribute} \emph{beer} facts and
    the \code{beer-why} template.
\end{enumerate}


\section{Objects}

The conceptualization phase brought several objects to light. In
Table~\ref{tab:objects} we report them as the \emph{deftemplates} through which
they are represented in working memory, describing what each represents and
listing its real slots together with their type, so as not to weigh down the
discussion with the enumeration of every individual fact.

\begin{longtable}{p{0.20\linewidth} p{0.42\linewidth} p{0.30\linewidth}}
\caption{Objects identified during the conceptualization phase, with their
real slots.}\label{tab:objects}\\
\toprule
\textbf{Object} & \textbf{Description} & \textbf{Attributes (slots / type)} \\
\midrule
\endfirsthead
\toprule
\textbf{Object} & \textbf{Description} & \textbf{Attributes (slots / type)} \\
\midrule
\endhead
\midrule
\multicolumn{3}{r}{\footnotesize\itshape continued on next page}\\
\endfoot
\bottomrule
\endlastfoot

\code{beer} &
A named beer belonging to one of fifteen recognised styles, described by the
symbolic descriptors over which the qualitative reasoning matches. &
\code{style} (STRING, allowed-strings); \code{name} (STRING);
\code{alcohol}, \code{color}, \code{flavor}, \code{fermentation},
\code{carbonation} (multislot SYMBOL, each with allowed-symbols);
\code{link} (STRING). \\
\addlinespace

\code{beer-profile} &
The authoritative per-style sensory profile extracted from the CraftBeer~2017
guide: descriptors, palate, hop/malt character and serving advice. &
\code{name} (STRING); \code{alcohol}, \code{colour}, \code{clarity},
\code{body}, \code{palate-carbonation}, \code{finish}, \code{hop}, \code{malt},
\code{glass}, \code{temperature}, \code{country} (STRING). \\
\addlinespace

\code{beer-range} &
The per-style numeric ranges extracted from the same guide; these data ground
the fuzzy membership functions on the measurable axes. &
\code{name} (STRING); \code{abv}, \code{ibu}, \code{srm}, \code{og},
\code{fg}, \code{co2} (multislot). \\
\addlinespace

\code{pairing} &
The authoritative food pairing for a style (cheese, entr\'ee and dessert),
matched against the foods named by the user. &
\code{style} (STRING); \code{cheese}, \code{entree}, \code{dessert} (STRING). \\
\addlinespace

\code{attribute} &
A graded piece of evidence about the desired profile, in the canonical
\emph{best-*} form (e.g. \code{best-color}, \code{best-alcohol},
\code{best-flavor}, \code{best-style}, \code{best-name}), and the final
recommendation facts (\code{name beer}). &
\code{name}; \code{value}; \code{certainty} (FLOAT, range $[-1,1]$,
default $1.0$). \\
\addlinespace

\code{ling-var} &
A linguistic variable: one measurable axis together with the bounds of its
universe of discourse. &
\code{axis}; \code{from}; \code{to}. \\
\addlinespace

\code{ling-term} &
A linguistic term of an axis, represented as a trapezoidal membership function
$a\le b\le c\le d$ (shoulders when $a{=}b$ or $c{=}d$), with breakpoints derived
from the data distributions. &
\code{axis}; \code{term}; \code{a}, \code{b}, \code{c}, \code{d}. \\
\addlinespace

\code{desire} &
A weighted preference for a linguistic term on a continuous axis, derived from
the \emph{best-*} evidence in fuzzy mode. &
\code{axis}; \code{term}; \code{weight} (default $1.0$). \\
\addlinespace

\code{desire-cat} &
A weighted preference for a value on a categorical (symbolic) slot of a beer,
such as \code{flavor} or \code{carbonation}. &
\code{slot}; \code{value}; \code{weight} (default $1.0$). \\
\addlinespace

\code{suggest-name} &
A direct, weighted suggestion of a specific beer name (e.g. produced by a meal
pairing). &
\code{name}; \code{weight} (default $1.0$). \\
\addlinespace

\code{suggest-style} &
A direct, weighted suggestion of a specific beer style. &
\code{style}; \code{weight} (default $1.0$). \\
\addlinespace

\code{UI-state} &
A single turn of the dialogue: the message shown to the user, its help/why
text, the admissible answers and the recorded response, together with the
relation it asserts and its position in the conversation. &
\code{id} (gensym default); \code{display}; \code{help}; \code{why};
\code{relation-asserted} (default \code{none}); \code{valid-answers}
(multislot); \code{response} (default \code{none}); \code{state}
(default \code{middle}); \code{translated} (default \code{FALSE}). \\
\addlinespace

\code{state-list} &
The bookkeeping of the conversation: the current turn and the sequence of turns
visited, enabling navigation and clean-up. &
\code{current}; \code{sequence} (multislot). \\
\addlinespace

\code{translation} &
A gettext-style localisation overlay: the English string is the source key, the
Italian string is the optional overlay (with English fallback). &
\code{en} (STRING); \code{it} (STRING). \\
\addlinespace

\code{meal-keyword} &
A food the user mentioned during the dialogue, matched against the authoritative
\code{pairing} facts. &
\code{text}. \\
\addlinespace

\code{beer-why} &
The per-beer rationale produced by the fuzzy backend (which factors matched and
how strongly), used to build a data-grounded explanation. &
\code{name}; \code{text}. \\
\end{longtable}


\section{Relations and Properties}

Table~\ref{tab:relations} reports the relations among the objects and the
properties that link them. The first group describes how the data-grounded
descriptions qualify a beer or style; the second how evidence is accumulated and
turned into desires; the third how the dialogue is modelled and localised.

\begin{longtable}{p{0.46\linewidth} p{0.46\linewidth}}
\caption{Relations and properties of the objects.}\label{tab:relations}\\
\toprule
\textbf{Relation} & \textbf{Description} \\
\midrule
\endfirsthead
\toprule
\textbf{Relation} & \textbf{Description} \\
\midrule
\endhead
\midrule
\multicolumn{2}{r}{\footnotesize\itshape continued on next page}\\
\endfoot
\bottomrule
\endlastfoot

A \code{beer} \emph{belongs to} a \code{style} &
Each beer carries a \code{style} drawn from the fifteen recognised styles and is
described by its symbolic \emph{axes} (colour, alcohol, flavour, fermentation,
carbonation), over which both backends match candidates. \\
\addlinespace

A \code{beer-profile} \emph{describes} a beer &
Joined by \code{name}, the profile supplies the authoritative sensory
description and the serving hints (glass, temperature) appended to a
recommendation. \\
\addlinespace

A \code{beer-range} \emph{measures} a beer &
Joined by \code{name}, the numeric ranges (ABV, IBU, SRM, \dots) provide the
crisp value---the midpoint of the range---fed to the fuzzy membership functions
on the colour, alcohol and bitterness axes. \\
\addlinespace

A \code{pairing} \emph{recommends} a style &
Joined by \code{style}, the authoritative cheese/entr\'ee/dessert pairing is
matched against the \code{meal-keyword} facts to grant a pairing bonus. \\
\addlinespace

An \code{attribute} \emph{carries best-* evidence with a certainty} &
Each \emph{best-*} attribute records a desired value on one axis (or a direct
style/name suggestion) together with its certainty in $[-1,1]$; rules that assert
the same value on the same axis are merged into a single accumulated certainty. \\
\addlinespace

A \code{ling-var} \emph{spans} an axis; a \code{ling-term} \emph{partitions} it &
Each measurable axis (colour~$\leftarrow$~SRM, alcohol~$\leftarrow$~ABV\%,
bitterness~$\leftarrow$~IBU) has a universe of discourse and is covered by
overlapping trapezoidal terms whose breakpoints are derived from the data
distributions. \\
\addlinespace

A \code{desire} / \code{desire-cat} \emph{weights} a preference &
In fuzzy mode the \emph{best-*} evidence is translated into desires: a
\code{desire} weights a linguistic term on a continuous axis, while a
\code{desire-cat} weights a value on a categorical slot; both carry a weight in
$[-1,1]$. \\
\addlinespace

A \code{suggest-name} / \code{suggest-style} \emph{points to} a candidate &
A direct, weighted suggestion of a specific beer name or style, typically
produced by a meal pairing, contributing to that candidate's score. \\
\addlinespace

A \code{beer-why} \emph{explains} a recommendation &
Joined by \code{name}, it records which desires the beer satisfied and how
strongly, so that the final explanation is data-grounded. \\
\addlinespace

A \code{UI-state} \emph{asserts} a relation and \emph{advances} the dialogue &
Each turn shows a message, offers \code{valid-answers}, records the
\code{response} and asserts the corresponding fact; its \code{state}
(initial / middle / final) marks the phase of the conversation. \\
\addlinespace

A \code{state-list} \emph{orders} the conversation &
It keeps the \code{current} turn and the \code{sequence} of visited turns,
enabling navigation through the dialogue and final clean-up. \\
\addlinespace

A \code{translation} \emph{overlays} a label &
Keyed on the English source string, it provides the localised (Italian) text,
falling back to English when no overlay exists. \\
\addlinespace

A \code{meal-keyword} \emph{matches} a pairing &
Each food named by the user is compared against the cheese/entr\'ee/dessert text
of the \code{pairing} facts to reward styles whose authoritative pairing matches
the user's meal. \\
\end{longtable}


\section{The two reasoning models (overview)}

Both backends---selected by the \code{?*backend*} global (\code{cf} or
\code{fuzzy})---consume the same \emph{best-*} evidence accumulated on the
colour, alcohol, flavour, fermentation and carbonation axes during the dialogue.
They differ only in how that evidence is combined and matched against the
candidate beers. We give here a conceptual overview; the formal calculus is
detailed in the Implementation chapter.

\subsection{Certainty factors}

The first model is the classical MYCIN scheme. Every piece of evidence is a
graded judgement carrying a certainty factor in $[-1,1]$, ranging from full
disbelief through ignorance to full belief. When several rules independently
assert the same value on the same axis, their certainties are merged by the MYCIN
combination function, which accumulates concordant evidence towards the extremes
while letting discordant evidence cancel. A beer is then scored by accumulating
the certainties of the axes whose symbolic value it matches, and the highest-%
scoring candidates are returned. The model is transparent and faithful to expert
practice, but the weights are assigned by the expert and are not quantified by
the source guides; this intrinsic arbitrariness, and the coarse ties it produces,
is precisely what motivates the fuzzy alternative.

\subsection{Fuzzy logic}

The second model replaces the hand-assigned certainties on the \emph{measurable}
axes with objective, data-grounded memberships. The measurable axes are modelled
as \emph{linguistic variables}---colour over SRM, alcohol over ABV\% and a new
bitterness axis over IBU---each covered by overlapping linguistic terms encoded
as trapezoidal membership functions. Crucially, the breakpoints are
\emph{derived from data}: crossing the symbolic labels of the dataset with the
numeric ranges of the guide yields the empirical distribution of each term
(validated against the p25 / median / p75 quantiles), from which the trapezoids
are built. The breakpoints currently in use are summarised in
Table~\ref{tab:trapezoids} as a forward reference.

\begin{table}[ht]
\centering\small
\caption{Data-grounded trapezoidal breakpoints of the linguistic terms (forward
reference; derivation in the Implementation chapter).}\label{tab:trapezoids}
\begin{tabular}{llcccc}
\toprule
axis & source & \multicolumn{4}{c}{terms (trapezoid breakpoints $a,b,c,d$)}\\
\midrule
colour & SRM & pale\,(0,0,6,9) & amber\,(6,9,11,14) & brown\,(11,14,21,27) & dark\,(21,27,50,50)\\
alcohol & ABV\% & n.d.\,(2,2,4.5,5.2) & mild\,(4.5,5,6,6.6) & noticeable\,(6,6.6,8,8.6) & harsh\,(8,8.6,15,15)\\
bitterness & IBU & low\,(0,0,20,25) & medium\,(20,25,30,35) & high\,(30,35,70,70) &\\
\bottomrule
\end{tabular}
\end{table}

Carbonation, by contrast, stays symbolic: the guide's CO\textsubscript{2} volumes
do not separate \emph{medium} from \emph{high}, so a data-driven membership
function would not be justified. In fuzzy mode the \emph{best-*} evidence is
translated into \emph{desires}---weighted preferences on the linguistic axes---
categorical flavour and carbonation desires, and direct style/name suggestions
arising from the meal. Each beer is then ranked by a normalised weighted-mean
satisfaction score: the membership of the beer's measured value in each desired
term, weighted by the desire and normalised by the total weight, with a beer
missing data on a desired axis penalised rather than exempted. A further term
boosts styles whose authoritative CraftBeer pairing matches a food named by the
user. Here the gain is real: the memberships are objective, so the fuzzy ranking
grades the candidates finely where the certainty-factor model can only produce
coarse ties. Membership evaluation, the fuzzy set operations and the centroid
defuzzification are provided by a reusable, agnostic pure-CLIPS fuzzy library; the
corresponding equations and the equivalence with the original FuzzyCLIPS
primitives are given in the Implementation chapter.
